\chapter{Mathematical Structure of Reduction and Continuum Claims}
\index{reduction!mathematical structure|textbf}

\label{ch:mathematics-of-reduction}

\chapterstatus{proof contracts and proof strategies are presented with explicit status; no proposed physics proof package is represented as complete or reviewed.}


\section{Excluded-space reduction}
\index{excluded-space reduction}
\index{downfolding}


Let $P+Q=I$ be complementary orthogonal projectors and write the block
components of a self-adjoint operator $H$ as $H_{PP}=PHP$,
$H_{PQ}=PHQ$, and so on. When a real energy $E$ belongs to the resolvent set of
$H_{QQ}$, eliminating the $Q$ component gives the energy-dependent Feshbach
effective operator~\cite{feshbach1958,feshbach1962}
\begin{equation}
  H_{\mathrm{eff}}(E)
  =
  H_{PP}+H_{PQ}(E-H_{QQ})^{-1}H_{QP}.
\end{equation}
If
\begin{equation}
  \Delta_Q=\operatorname{dist}(E,\sigma(H_{QQ}))>0,
\end{equation}
then the self-adjoint resolvent estimate suggests the proposed bound
\begin{equation}
  \left\|
    H_{PQ}(E-H_{QQ})^{-1}H_{QP}
  \right\|
  \leq
  \frac{\|H_{PQ}\|\,\|H_{QP}\|}{\Delta_Q}.
\end{equation}
For $H_{QP}=H_{PQ}^{\dagger}$ this becomes
$\|H_{PQ}\|^2/\Delta_Q$.

The \texttt{PRF-20.02} obligation must keep this real off-spectrum setting
separate from open-system or analytically continued resonance formulations,
where non-Hermiticity and widths require additional theory. No small correction
is inferred without both a coupling estimate and resolvent distance.

\section{Crossover and operator-to-observable bounds}
\index{operator-to-observable bound}
\index{crossover!continuum}

A continuum comparison may contain several mathematically distinct limiting
operations. A continuum-discretization limit refines the numerical mesh for a
fixed continuum operator; a discrete-to-continuum limit changes the represented
lattice scale and requires an explicit identification between changing state
spaces~\cite{nakamuraTadano2021}; a band-edge homogenization limit weakens and
widens a defect relative to the crystal period while deriving an envelope
operator~\cite{hoeferWeinstein2011,ducheneVukicevicWeinstein2015}; and an
infinite-volume limit enlarges the domain while holding local physics fixed.
Changing the defect profile or switching between fixed-peak and
fixed-integrated normalization is a model-family change unless the scaling is
part of the declared limit. Convergence on a path that varies several of these
quantities therefore does not identify which approximation produced the
observed agreement.

The sources also differ in operator class and spectral claim. Weak-binding
results concern existence below a threshold or in periodic gaps
~\cite{simon1976weak,parzygnat2010}; homogenization results concern modes
bifurcating from a band edge under weak/slow scaling; and finite-periodic-volume
formulas concern image shifts under separate range and tail assumptions
~\cite{koenigLeeHammer2011}. None implies that a finite represented
multiorbital operator satisfies another source's hypotheses. Numerical
agreement of one energy or projector therefore cannot substitute for the
state-space map, operator residual, and scaling assumptions of the applicable
result.

Let $D_d=\Delta H_d-V_{\mathrm{cont},d}$ be an aligned atomistic-minus-continuum
discrepancy in a common space, and let $P_{>R}$ project onto a declared exterior
region. Define
\begin{equation}
  \eta_d(R)
  =
  \|P_{>R}D_dP_{>R}\|.
\end{equation}
When the exterior subspaces are nested, compression by a smaller exterior
projector cannot increase the operator norm. A crossover criterion may then be
defined by
\begin{equation}
  r_{c,d}(\tau)
  =
  \inf\{R:\eta_d(R)\leq\tau\}.
\end{equation}

The \texttt{PRF-30.01} obligation must distinguish nonemptiness of the set,
existence of the infimum, attainment by a tested radius, numerical estimation,
and physical interpretation. The assumption
$\eta_d(R)\to0$ as $R\to\infty$ is an asymptotic-locality hypothesis, not a
consequence established by general Kohn--Sham DFT.

\citationtodo{high priority}{Check the asymptotic-locality discussion against
nearsightedness literature while preserving the statement as an explicit
hypothesis. Candidate: Prodan and Kohn (2005), provisional key
\texttt{prodanKohn2005}; the candidate record is present for this editorial note.}

For self-adjoint atomistic and reduced operators with perturbation
$E=H_{\mathrm{atom}}-H_{\mathrm{red}}$, standard perturbation theory motivates
spectral-distance bounds controlled by $\|E\|$. Invariant-subspace rotation can
be controlled by a Davis--Kahan-type estimate when the target cluster is
separated by a declared spectral gap~\cite{davis1970}. The proposed
\texttt{PRF-30.02} result must identify the finite or infinite setting, norm,
spectral window, gap, state identification, and degeneracy treatment. A
worst-case bound is not assumed sharp, and a small spectral error does not imply
a small operator error.

\citationtodo{medium priority}{Check the general spectral-perturbation statement
against Kato's standard treatment, provisional key \texttt{kato1995}, while
retaining \texttt{davis1970} for the named
Davis--Kahan invariant-subspace estimate.}

\section{Continuum fitting and identifiability}
\index{identifiability!continuum}


For continuum parameters $\boldsymbol\theta\in\Theta$, define a declared
exterior objective
\begin{equation}
  J_R(\boldsymbol\theta)
  =
  \left\|
    P_{>R}
    [\Delta H_d-V_{\mathrm{cont}}(\boldsymbol\theta)]
    P_{>R}
  \right\|.
\end{equation}
If $\Theta$ is nonempty and compact and $J_R$ is continuous, the extreme-value
theorem gives existence of a minimizer. The proposed \texttt{PRF-30.03}
program separates this existence result from uniqueness, conditioning,
stability, and physical identifiability. An optimizer output supplies none of
those stronger properties automatically.

\section{Compatibility and certified incompatibility}
\index{compatibility!certified incompatibility}


Let $\mathfrak A_E$ and $\mathfrak A_H$ be nonempty spectral- and
operator-admissible subsets of a declared metric model space. If both sets are
compact, their distance
\begin{equation}
  \delta^{\ast}
  =
  \min_{H_E\in\mathfrak A_E,\,H_H\in\mathfrak A_H}
  d(H_E,H_H)
\end{equation}
is attained. If the compact sets are disjoint, the distance is positive. This
standard existence argument is only part of \texttt{PRF-40.01}: the scientific
application must also establish closure or compactness, continuity of the
losses, and a reproducible admissible gauge group.

A failed local optimizer does not prove that the sets are disjoint or that
$\delta^{\ast}>0$. The proposed \texttt{PRF-40.03} certification obligation
requires a valid global lower bound, for example from analytic estimates,
interval methods, branch-and-bound, exhaustive certified reduction, or an
applicable validated convex relaxation. \texttt{PRF-40.04} then asks what can be
inferred when a nested model class reduces the residual; a smaller residual does
not uniquely identify the physical cause of the improvement.

\citationtodo{medium priority}{Check the certified global-optimization examples,
especially interval and branch-and-bound methods, against a validated global
optimization source. Candidate: Neumaier (2004), provisional key
\texttt{neumaier2004}; the candidate record is present for this editorial note.}

\section{Reduction-path commutativity}
\index{commutativity!reduction paths}


For aligned operands in one common space, let $\mathcal R$ be a reduction map.
The path defect is
\begin{equation}
  \varepsilon_{\mathrm{path}}
  =
  \left\|
    \mathcal R(H_d-H_b^{\mathrm{aligned}})
    -
    [\mathcal R(H_d)-\mathcal R(H_b^{\mathrm{aligned}})]
  \right\|.
\end{equation}
If the same linear map acts on both operands in the same coordinates, exact
commutativity follows from linearity. This elementary fact does not apply
automatically to nonlinear fitting, separately selected model classes,
gauge-dependent truncation, or independently chosen alignment maps.
Appendix~\ref{app:one-dimensional-reduction} supplies a related controlled route
comparison: matched direct and mediated Fourier projections agree, whereas
altered reciprocal weights or fitting domains produce nonzero route defects.
That experiment does not instantiate the impurity-difference expression above;
it demonstrates the same need to hold maps, coordinates, and objectives fixed
before interpreting path agreement.

The \texttt{PRF-20.04} program distinguishes gauge equivariance from path
commutativity and seeks bounds when exact equality fails. The frozen
\texttt{PRF-05.06} contract addresses gauge invariance of a residual only after
the two paths are already conformant.

This conformity requirement is operational rather than merely notational.
Localized-basis disorder and electron--defect methods require a common Wannier
representation or explicit Bloch--Wannier maps before defect contributions are
transported~\cite{berlijnVoljaKu2011,luParkZhouBernardi2020}. Model-stitching
work aligns orbital pairing, spin orientation, phase, and scalar potential
before minimizing overlap-region Hamiltonian differences~\cite{lihmPark2019}.
A recent preprint similarly uses assignment, a dense unitary Procrustes map, and
an induced overlap metric when embedding independently generated clean and
defect Wannier models~\cite{shiXieZhangChuGong2026}. These constructions are
precedents for explicit identification; they do not prove that an inferred map
is unique, well conditioned, or complete on an unobserved sector.
